% GENERATED FILE. Edit canonical lessons, then run npm run generate:books.
% canonical-source-hash: fnv1a64:7435b280da5bdbbc
% canonical-lessons: FA-C298-agar, FA-C298-zhakat, FA-W298-zhe, FA-C298-agar-baran-biyayad, FA-C298-bastegi-darad, FA-R298-if-and-it-depends

\chapter{If It Rains}
\label{ch:fa-if-it-rains}
\hlchaptermodality{298}

\begin{chapteropening}
\textbf{By the end of this chapter:} \emph{I can state a real condition with agar and the subjunctive, give its result, and answer a plan with bastegi darad.}

Complete FA-R298-if-and-it-depends: a plan that depends on the weather, said with agar and bastegi darad.
\end{chapteropening}

\section[\texorpdfstring{\fa{اگر} … \fa{باشد} (agar ... bâshad)}{اگر … باشد (agar ... bâshad)}]{\fa{اگر} … \fa{باشد} (agar ... bâshad) — if it is}

\label{lesson:FA-C298-agar}



\begin{quote}
Say \textquotedblleft{}if\textquotedblright{}, then \textquotedblleft{}otherwise\textquotedblright{}. (\emph{Agar}, \emph{vagarna}.) Then the subjunctive of \emph{budan} for \textquotedblleft{}she\textquotedblright{}: \textquotedblleft{}that she be\textquotedblright{}. (\emph{Bâshad.})

\emph{Agar} has waited since chapter one hundred and fifty-eight. Now you hold the verb shape it asks for.
\end{quote}

\begin{grammarlens}[title={Grammar Lens: the condition first, in the subjunctive}]
\textbf{\fa{اگر}}, \emph{agar}, opens the condition. The condition usually comes first, then a pause, then what follows from it:

\begin{quote}
\textbf{\fa{اگر} \fa{هوا} \fa{خوب} \fa{باشد،} \fa{به} \fa{پارک} \fa{میرویم}.} — \emph{agar havâ khub bâshad, be pârk miravim} — If the weather is good, we'll go to the park.
\end{quote}

The condition is about later, and it may not come true, so its verb is in the subjunctive: \emph{bâshad}, not \emph{ast}. The second half is a plain present, which in Persian also covers the future.

\begin{quote}
\textbf{\fa{اگر} \fa{خسته} \fa{باشی،} \fa{بهتر} \fa{است} \fa{بخوابی}.} — \emph{agar khaste bâshi, behtar ast bekhâbi} — If you're tired, you'd better sleep.
\end{quote}

\emph{Agar} can also sit after the person it is about: \emph{to agar khaste bâshi...}. Either way, \emph{agar} stays inside the condition half.
\end{grammarlens}

\begin{grammarlens}[title={Grammar Lens: vagarna, the other branch}]
\emph{Vagarna}, \textquotedblleft{}otherwise\textquotedblright{}, is the condition turned round, \textquotedblleft{}and if not\textquotedblright{}:

\begin{quote}
\textbf{\fa{زود} \fa{بروید،} \fa{وگرنه} \fa{قطار} \fa{را} \fa{از} \fa{دست} \fa{میدهید}.} — \emph{zud beravid, vagarna qatâr râ az dast midahid} — Leave early, otherwise you'll miss the train.
\end{quote}
\end{grammarlens}

\begin{cousinweb}
\emph{agar} continues Middle Persian \emph{agar}, \textquotedblleft{}if\textquotedblright{} (EtymDB 2.1). \emph{Vagarna} is \emph{va}, \textquotedblleft{}and\textquotedblright{}, \emph{agar}, and \emph{na}, \textquotedblleft{}not\textquotedblright{}, run together: the same \textquotedblleft{}if\textquotedblright{} with \textquotedblleft{}not\textquotedblright{} on the end.
\end{cousinweb}

\subsection*{Your turn}
Say these aloud:

\begin{itemize}
\raggedright
  \item \emph{agar havâ khub bâshad, be pârk miravim}
  \item \emph{agar khaste bâshi, behtar ast bekhâbi}
  \item \emph{zud beravid, vagarna qatâr râ az dast midahid}
  \item a plan of your own for tomorrow, with \emph{agar} and \emph{bâshad}
\end{itemize}

\begin{tcolorbox}[breakable,colback=teal!4,colframe=teal!35!black,title={Before you move on}]
Which half comes first? (Usually \textbf{the condition}, with \emph{agar}.) Why \emph{bâshad} and not \emph{ast}? (The condition is about later, and \textbf{may not come true}.) \textquotedblleft{}Otherwise\textquotedblright{}? (\emph{\textbf{Vagarna.}})
\end{tcolorbox}
\section[\texorpdfstring{\fa{ژاکت} (zhâkat)}{ژاکت (zhâkat)}]{\fa{ژاکت} (zhâkat) — a cardigan, in case it turns cold}

\label{lesson:FA-C298-zhakat}



\begin{quote}
Say \textquotedblleft{}if the weather is good\textquotedblright{}. (\emph{Agar havâ khub bâshad.}) Then a padded jacket and a pullover. (\emph{Kâpshan}, \emph{poliver}.)

A condition is easiest to practise on the weather, and on what to wear for it.
\end{quote}

\subsection*{You'll want to know: \fa{ژاکت}}
\textbf{\fa{ژاکت}} — \emph{zhâkat} — \textquotedblleft{}a cardigan\textquotedblright{}, a knitted jacket that buttons down the front.

\begin{quote}
\textbf{\fa{اگر} \fa{هوا} \fa{سرد} \fa{باشد،} \fa{ژاکت} \fa{بپوشید}.} — \emph{agar havâ sard bâshad, zhâkat bepushid} — If it's cold, put on a cardigan.
\end{quote}

\emph{Bepushid} is \emph{pushidan}, \textquotedblleft{}to put on, to wear\textquotedblright{}, on its stem \emph{push-}. A \emph{zhâkat} is lighter than a \emph{kâpshan}, the padded jacket from chapter one hundred and twenty, and it opens at the front where a \emph{poliver} goes over the head.

Say \emph{zhâ-kat}. The \emph{zh} is the sound in the middle of English \emph{measure}, and its letter is one Persian added to the Arabic alphabet. The next lesson puts it in your hand.

\begin{cousinweb}
EtymDB 2.1 has no entry for \emph{zhâkat}. It is usually traced to French \emph{jaquette}, a jacket. EtymDB takes English \textbf{jacket} back to Middle French \emph{jacquet}, from the same French family.
\end{cousinweb}

\subsection*{Your turn}
Say these aloud:

\begin{itemize}
\raggedright
  \item \emph{zhâkat}, with the \emph{zh} of \emph{measure}
  \item \emph{agar havâ sard bâshad, zhâkat bepushid}
  \item a comparison — \emph{zhâkat az kâpshan saboktar ast}, a cardigan is lighter than a padded jacket
  \item the same advice to a friend — \emph{zhâkat bepush}
\end{itemize}

\begin{tcolorbox}[breakable,colback=teal!4,colframe=teal!35!black,title={Before you move on}]
\textquotedblleft{}A cardigan\textquotedblright{}? (\emph{\textbf{Zhâkat.}}) \textquotedblleft{}If it's cold, put on a cardigan\textquotedblright{}? (\emph{\textbf{Agar havâ sard bâshad, zhâkat bepushid.}}) Which English word comes from the same French family? (\textbf{Jacket.})
\end{tcolorbox}
\section[\texorpdfstring{\fa{ژ} (zhe)}{ژ (zhe)}]{\fa{ژ} (zhe) — the letter that opens \fa{ژاکت}}

\label{lesson:FA-W298-zhe}



\begin{quote}
Before the new one: write \fa{ز} and \fa{پ} once each, and name them.

Now say \textbf{\fa{ژاکت}} (\emph{zhâkat}), \textquotedblleft{}a cardigan\textquotedblright{}. You met it a moment ago. This lesson takes its first letter out and puts it in your hand.
\end{quote}

\subsection*{Script you'll notice: \fa{ژ}}
\textbf{\fa{ژ}} — \emph{zhe}, the sound in the middle of English \emph{measure}.

It opens \textbf{\fa{ژاکت}}. It is \textbf{\fa{ز}} with \textbf{three dots} in place of one, and like \textbf{\fa{ز}} it takes nothing on its left: after \textbf{\fa{ژ}} the pen lifts, and \textbf{\fa{ا}} stands on its own before \textbf{\fa{کت}}.

Arabic has no \emph{zh}, so the Arabic alphabet had no letter for it. Persian made four letters of its own from Arabic ones, with extra dots or an extra stroke: \textbf{\fa{پ}} from \textbf{\fa{ب}}, \textbf{\fa{چ}} from \textbf{\fa{ج}}, \textbf{\fa{گ}} from \textbf{\fa{ک}}, and \textbf{\fa{ژ}} from \textbf{\fa{ز}}. You have now written all four.

\subsection*{Writing: \fa{ژ}}
Put your pen on \fa{ژ} and follow its line. Copy the shape you can see — slowly, and larger than it is printed.

\begin{quote}
This book does not yet tell you \textbf{where to start this letter or which way to travel}. It is not written down here until it can be written down with a source. Copying what is in front of you needs no such source, and it is how the shape gets into your hand in the meantime.
\end{quote}

\subsection*{Your turn}
\begin{itemize}
\raggedright
  \item \emph{Find:} \fa{ژ} at the start of \textbf{\fa{ژاکت}}, and point at its three dots
  \item \emph{Trace it:} \fa{ژ} three times, saying \emph{zh} as you finish each one
  \item \emph{Write it:} \fa{ز}, then \fa{ژ}, from memory: one dot, then three
  \item \emph{Write it:} \textbf{\fa{ژاکت}}, lifting the pen after \fa{ژ} and after \fa{ا}
  \item \emph{Read it:} \textbf{\fa{ژاکت}} aloud once more
\end{itemize}

\begin{tcolorbox}[breakable,colback=teal!4,colframe=teal!35!black,title={Before you move on}]
Which letter is this — \fa{ژ}? (\emph{\textbf{zhe}}, the \emph{zh} of \emph{measure}.) Which word did it come from? (\textbf{\fa{ژاکت}}, \emph{zhâkat}, \textquotedblleft{}a cardigan\textquotedblright{}.) Name the four letters Persian added. (\textbf{\fa{پ} \fa{چ} \fa{ژ} \fa{گ}}.)
\end{tcolorbox}
\section[\texorpdfstring{\fa{اگر} \fa{باران} \fa{بیاید} (agar bârân biyâyad)}{اگر باران بیاید (agar bârân biyâyad)}]{\fa{اگر} \fa{باران} \fa{بیاید} (agar bârân biyâyad) — if it rains}

\label{lesson:FA-C298-agar-baran-biyayad}



\begin{quote}
Say \textquotedblleft{}if it's cold, put on a cardigan\textquotedblright{}. (\emph{Agar havâ sard bâshad, zhâkat bepushid.}) Then \textquotedblleft{}to come\textquotedblright{}, and its present stem. (\emph{Âmadan}, \emph{â-}.)

In everyday Persian, rain comes: \emph{bârân miâyad}, \textquotedblleft{}it's raining\textquotedblright{}, is word for word \textquotedblleft{}rain is coming\textquotedblright{}. That puts \emph{âmadan} into a rainy condition.
\end{quote}

\begin{grammarlens}[title={Grammar Lens: biyâyad, the subjunctive of âmadan}]
\emph{Be-} meets the stem \emph{â-}, and a \emph{y} slips in to keep the two vowels apart: \textbf{\fa{بیاید}}, \emph{biyâyad}, \textquotedblleft{}that it come\textquotedblright{}. So \emph{biyâyam}, \emph{biyâyi}, \emph{biyâyad}.

\begin{quote}
\textbf{\fa{اگر} \fa{باران} \fa{بیاید،} \fa{در} \fa{خانه} \fa{میمانیم}.} — \emph{agar bârân biyâyad, dar khâne mimânim} — If it rains, we'll stay at home.
\end{quote}

Word for word, \textquotedblleft{}if rain come, we stay in the house\textquotedblright{}. \emph{Mimânim} is \emph{mândan}, \textquotedblleft{}to stay\textquotedblright{}, on its stem \emph{mân-}.
\end{grammarlens}

\begin{grammarlens}[title={Grammar Lens: three shapes for the result}]
A real condition can come true, and its result takes one of three shapes.

What will happen, in the plain present:

\begin{quote}
\textbf{\fa{اگر} \fa{دیر} \fa{برسی،} \fa{قطار} \fa{را} \fa{از} \fa{دست} \fa{میدهی}.} — \emph{agar dir beresi, qatâr râ az dast midahi} — If you arrive late, you'll miss the train.
\end{quote}

What to do, as a polite instruction: \emph{agar bârân biyâyad, zhâkat bepushid}, if it rains, put on a cardigan.

What must happen, with \emph{bâyad}: \emph{agar bârân biyâyad, bâyad zud beravim}, if it rains, we have to leave early.

The \emph{agar} half keeps its subjunctive in all three. Only the result changes.
\end{grammarlens}

\subsection*{Your turn}
Say these aloud:

\begin{itemize}
\raggedright
  \item \emph{biyâyam, biyâyi, biyâyad}
  \item \emph{agar bârân biyâyad, dar khâne mimânim}
  \item \emph{agar dir beresi, qatâr râ az dast midahi}
  \item one condition, three results — \emph{agar bârân biyâyad}, then a present, an instruction, and \emph{bâyad}
\end{itemize}

\begin{tcolorbox}[breakable,colback=teal!4,colframe=teal!35!black,title={Before you move on}]
\textquotedblleft{}That it come\textquotedblright{}? (\emph{\textbf{Biyâyad.}}) Why the \emph{y}? (To keep \textbf{two vowels apart}.) What are the three shapes of the result? (\textbf{A present, a polite instruction, \emph{bâyad}}.) What stays the same in the \emph{agar} half? (\textbf{The subjunctive.})
\end{tcolorbox}
\section[\texorpdfstring{\fa{بستگی} \fa{دارد} (bastegi dârad)}{بستگی دارد (bastegi dârad)}]{\fa{بستگی} \fa{دارد} (bastegi dârad) — it depends}

\label{lesson:FA-C298-bastegi-darad}



\begin{quote}
Say \textquotedblleft{}if it rains, we'll stay at home\textquotedblright{}. (\emph{Agar bârân biyâyad, dar khâne mimânim.}) Then \textquotedblleft{}to tie, to close\textquotedblright{}. (\emph{Bastan.})

Asked about a plan, you often cannot answer yes or no yet. The honest answer is \textquotedblleft{}it depends\textquotedblright{}, and then the condition.
\end{quote}

\subsection*{You'll want to know: \fa{بستگی} \fa{دارد}}
\textbf{\fa{بستگی} \fa{دارد}} — \emph{bastegi dârad} — \textquotedblleft{}it depends\textquotedblright{}.

\begin{quote}
— \textbf{\fa{فردا} \fa{به} \fa{پارک} \fa{میروید؟}} — \emph{fardâ be pârk miravid?} — Are you going to the park tomorrow?
\end{quote}

\begin{quote}
— \textbf{\fa{بستگی} \fa{دارد}. \fa{اگر} \fa{هوا} \fa{خوب} \fa{باشد،} \fa{بله}.} — \emph{bastegi dârad. agar havâ khub bâshad, bale.} — It depends. If the weather is good, yes.
\end{quote}

To name what it depends on, add \emph{be}, \textquotedblleft{}to\textquotedblright{}: \textbf{\fa{بستگی} \fa{دارد} \fa{به} \fa{هوا}}, \emph{bastegi dârad be havâ}, \textquotedblleft{}it depends on the weather\textquotedblright{}. The \emph{be} phrase can also come first: \emph{be havâ bastegi dârad}. And \emph{be shomâ bastegi dârad}, \textquotedblleft{}it depends on you\textquotedblright{}, is a polite \textquotedblleft{}it's up to you\textquotedblright{}.

The verb is \emph{dâshtan}, so it takes no \emph{mi-}: \emph{dârad}.

\begin{cousinweb}
\emph{bastegi} is \emph{baste}, \textquotedblleft{}bound, tied\textquotedblright{}, with \emph{-gi}, which makes a noun of it: \textquotedblleft{}boundness\textquotedblright{}. \emph{Baste} is the participle of \emph{bastan}, \textquotedblleft{}to tie, to close\textquotedblright{}, built the way \emph{gozashte} was built from \emph{gozashtan}. So \emph{bastegi dârad be havâ} is \textquotedblleft{}it has a binding to the weather\textquotedblright{}.

EtymDB 2.1 takes \emph{bastan} back to an Indo-European root meaning \textquotedblleft{}to bind, to tie\textquotedblright{}, the same root it gives for English \textbf{bind}.
\end{cousinweb}

\subsection*{Your turn}
\begin{itemize}
\raggedright
  \item \emph{Say it:} \emph{bastegi dârad. agar havâ khub bâshad, bale.}
  \item \emph{Say it:} \textquotedblleft{}it depends on the weather\textquotedblright{}, both ways — \emph{bastegi dârad be havâ}, \emph{be havâ bastegi dârad}
  \item \emph{Say it:} \textquotedblleft{}it's up to you\textquotedblright{} — \emph{be shomâ bastegi dârad}
  \item \emph{Answer:} \emph{fardâ be muze miravid?} — with \emph{bastegi dârad} and a condition of your own
\end{itemize}

\begin{tcolorbox}[breakable,colback=teal!4,colframe=teal!35!black,title={Before you move on}]
\textquotedblleft{}It depends on the weather\textquotedblright{}? (\emph{\textbf{Bastegi dârad be havâ.}}) What is \emph{bastegi}, taken apart? (\emph{\textbf{Baste}}, \textquotedblleft{}bound\textquotedblright{}, with \emph{\textbf{-gi}}: \textbf{boundness}.) Which English verb shares its root? (\textbf{Bind.})
\end{tcolorbox}
\section[(review)]{Review — a plan that depends on the weather}

\label{lesson:FA-R298-if-and-it-depends}



\begin{quote}
Say \textquotedblleft{}if the weather is good\textquotedblright{}, then \textquotedblleft{}if it rains\textquotedblright{}. (\emph{Agar havâ khub bâshad}, \emph{agar bârân biyâyad}.)
\end{quote}

\begin{grammarlens}[title={What you've built}]
Four lessons, and a plan can now hang on something. \emph{Agar} opens the condition and takes the subjunctive, \emph{bâshad} or \emph{biyâyad}. The result is a present, a polite instruction, or \emph{bâyad}. And when you cannot say yet, \emph{bastegi dârad} holds the answer open.

Two friends, making a plan:

\begin{quote}
— \emph{fardâ be pârk miravim?} — \emph{bastegi dârad be havâ. agar havâ khub bâshad, miravim. agar bârân biyâyad, dar khâne mimânim.} — \emph{agar sard bâshad?} — \emph{zhâkat bepush!}
\end{quote}
\end{grammarlens}

\subsection*{Your turn}
Say these aloud:

\begin{itemize}
\raggedright
  \item the four lines of the plan, playing both friends
  \item a plan with a comparison — \emph{agar havâ behtar bâshad, be pârk miravim}
  \item \emph{agar havâ mesl-e diruz bâshad, be behtarin pârk-e shahr miravim}
  \item a fact and a plan — \emph{do hafte ast ke bârân miâyad; agar fardâ biyâyad, dar khâne mimânam}
  \item a past point and a plan — \emph{se ruz pish be pârk raftim; agar havâ khub bâshad, dobâre miravim}
\end{itemize}

\begin{tcolorbox}[breakable,colback=teal!4,colframe=teal!35!black,title={Before you move on}]
\textquotedblleft{}It depends on the weather\textquotedblright{}? (\emph{\textbf{Bastegi dârad be havâ.}}) What verb shape follows \emph{agar}? (\textbf{The subjunctive.}) Which letter opens \textbf{\fa{ژاکت}}? (\textbf{\fa{ژ}}, \emph{zhe}.) And which sits in the middle of \textbf{\fa{مثل}}? (\textbf{\fa{ث}}, \emph{se}.)
\end{tcolorbox}
